\section{Why the Composition Does Not Add on Language}
\label{sec:mechanism}

Three accounts follow, labelled by evidential status: measured, or hypothesis
consistent with measurement.

\subsection{Routing divides the data, not the update budget}
\label{sec:dilution}

\textbf{[Ordering MEASURED; arithmetic MEASURED; mechanism HYPOTHESIS.]} The
obvious account of the ordering is capacity dilution: MoR's single block sees
100\% of tokens while MoE and MoRE split the same FFN budget across six experts,
so each expert is trained less. \textbf{That account is arithmetically false
here, and ruling it out is what makes the surviving hypothesis worth stating.}

The budgets are matched by construction, and the matching cancels. Each MoE/MoRE
expert FFN holds $2 \cdot 4 \cdot d^2 = 524{,}288$ weights and MoR's single FFN
holds $2 \cdot 24 \cdot d^2 = 3{,}145{,}728$, exactly six times as many
(\S\ref{sec:engine}). A run sees $404.2$M tokens. So:

\begin{center}
{\dispsize
\setlength{\tabcolsep}{9pt}\renewcommand{\arraystretch}{1.25}
\small
\begin{tabular}{lccc}
\toprule
per FFN parameter & MoE & MoR & MoRE \\
\midrule
tokens seen      & 128.5 & 128.5 & 128.5 \\
FFN applications & 128   & 865   & \textbf{894} \\
\bottomrule
\end{tabular}
}
\end{center}

Tokens per parameter is \emph{identical} across all three arms---the token split
and the parameter split are both by $E$, so they cancel exactly. And once
recursion is counted, applications per parameter slightly \emph{favours} MoRE
(894) over MoR (865), because MoRE runs marginally deeper. Whatever costs MoRE
0.12 nats against MoR, it is not that its parameters are trained less.

That comparison also explains the easy half of the ordering. MoRE gets
$894/128 \approx 7\times$ the FFN applications per parameter that MoE gets, which
is simply what recursion buys, and MoRE beats MoE by 0.311 nats. The recursive
multiply is real and it is the entire source of MoRE's advantage over MoE.

\textbf{What remains, once update count is excluded, is the \emph{diversity} of
what each parameter block sees.} MoR's block is exposed to the whole token
distribution; each MoE/MoRE expert is exposed only to the router-selected slice
that reaches it, which is both narrower and systematically biased, since the
router selects it. Structure shared across the whole distribution---the bulk of
what a language model must learn at this scale---can be learned once by MoR's
block but has to be re-learned independently inside all six experts. Six partial
copies of shared structure fit in the same parameter budget as one complete copy
only by making each copy worse.

This is a hypothesis, not an isolated effect, and we label it as one. Because the
budget equalises \emph{total} FFN width ($6 \times 4 = 1 \times 24$), MoR is
exactly the $E{=}1$ point of a matched-width expert sweep and MoRE the $E{=}6$
point, and the hypothesis predicts loss rising with $E$ at fixed width; two points
cannot isolate it, and the $E{=}1$ endpoint also has no router, gate or balance
term. The experiment that would is an intermediate point that keeps the router
($E{=}2$ at width 12, or $E{=}3$ at width 8). The
account earns its place because it is the only one left once the
arithmetic above excludes the update-count version; because it predicts the
stratified result---MoRE's advantage over MoE grows with token difficulty
(Appendix~\ref{sec:stratified}), where an incompletely learned shared structure should hurt
most; because MoRE is the arm still improving fastest when the budget ends
(\S\ref{sec:budget}), as slower re-learning would imply; and because it yields the
falsifiable scale prediction stated in the Limitations section.

\paragraph{The same block, two halting policies.}
On the language arm the halt head saturates: $97.3\%$ of tokens run to the depth
budget, and the learned policy is worth $0.060$ nats against a depth-matched
uniform policy---$4\%$ of MoR's margin on the same control. On the arithmetic
task the same block, with the same halt head, the same ponder term and the same
$285\times$ halt-to-task gradient ratio, settles near four of seven steps, and
its margin against that control matches MoR's on the mean ($0.842$ against
$0.837$). The halting mechanism is therefore not intrinsically inert: a saturated
policy is a property of the setting, not of the composition. Which of task,
objective and data-generating structure is responsible is not identified by our
design; the cross-task comparison supports the weaker claim, that a saturated
policy on one task does not license a conclusion about the mechanism.

\subsection{The two mechanisms compete for one budget}
\label{sec:competition}

\textbf{[MEASURED on a calibration grid; the canonical-corpus form is measured on
the balance axis.]} This is the paper's central claim: it explains why the
composition does not merely fail to add but subtracts.

The grid below varies the same halting weight and asks about the
\emph{direction} of the trade-off as the weight moves, not the level at which
depth settles: a level can fail to transfer while a direction holds. Whether this
one does is untested on the full corpus, so this subsection's halting claim
carries the strength of a one-seed development grid.

Depth and specialisation are not independent knobs in MoRE the way they are in MoR
and MoE separately. On the calibration grid, recovering adaptive depth by raising
the halting weight destroys expert differentiation:

\begin{center}
{\dispsize
\setlength{\tabcolsep}{4pt}\renewcommand{\arraystretch}{1.25}
\begin{tabular}{lccc}
\toprule
halting weight & mean depth & AMI & val loss \\
\midrule
0.001 (frozen) & 2.698 & \textbf{0.2851} & 5.2352 \\
0.107 & 1.476 & \textbf{0.0025} & 5.3346 \\
0.500 & 1.072 & 0.0207 & 5.3306 \\
\bottomrule
\end{tabular}
}
\end{center}

AMI of 0.0025 is statistically indistinguishable from the shuffled-label control.
The mechanism is direct: fewer recursion steps means fewer expert applications per
token, which means fewer opportunities for experts to differentiate. \textbf{Tune
toward adaptive depth and specialisation collapses; preserve specialisation and
depth saturates.} Both mechanisms draw on the same per-token compute budget, and at
this parameter and data budget neither gets enough.

The competition has a second face: the balance term is load-bearing and cannot be traded for depth either---zeroing it collapses the load onto one expert and drives AMI to zero. Specialisation is real against per-run permutation controls (MoRE AMI $0.1803$ vs control $0.0249$, $z=13.43$; MoE $0.1637$ vs $0.0364$) and the two arms are \emph{not} distinguishable, so recursion does not improve routing just as routing does not preserve adaptive depth. Appendix~\ref{app:specialisation} gives the controls.

\textbf{We also tested, and falsified, the most attractive rival account.} Subspace trapping---that routing once confines a token to one expert and repeated application drives the state to a fixed point---predicts lower dimensionality, oscillation and rescaling updates in the composed arm. All three fail: participation ratio is \emph{higher} in the recursive arms (57.3 and 70.3 against MoE 41.5), velocity decays monotonically, and updates are near-orthogonal (Appendix~\ref{app:trapping}). Excluding it leaves the budget competition as the account the data support.
